\section{Experiments}
% Centralized placeholders for the compact ablation and gradient audit.
% Replace each ``--'' with the measured value; the table and analysis text
% will update together.  Accuracy entries are macro means over the six
% benchmarks.  Cosines should be reported as mean \pm uncertainty.

% Component ablation: Macro Avg@8 / Macro Best@8.
\newcommand{\AblOPDAvg}{--}
\newcommand{\AblOPDBest}{--}
\newcommand{\AblOutcomeAvg}{--}
\newcommand{\AblOutcomeBest}{--}
\newcommand{\AblPrefixAvg}{--}
\newcommand{\AblPrefixBest}{--}
\newcommand{\AblOutcomePrefixAvg}{--}
\newcommand{\AblOutcomePrefixBest}{--}
\newcommand{\AblFullAvg}{--}
\newcommand{\AblFullBest}{--}

% Independent-reference gradient audit: cosine / negative-batch rate (%) /
% near-zero method-gradient rate (%) / relative norm.  AuditN is the
% common number of reference-valid rollout pairs used by every method.
\newcommand{\AuditN}{--}
\newcommand{\AuditOutcomeCos}{--}
\newcommand{\AuditOutcomeNeg}{--}
\newcommand{\AuditOutcomeZero}{--}
\newcommand{\AuditOutcomeNorm}{--}
\newcommand{\AuditRawCos}{--}
\newcommand{\AuditRawNeg}{--}
\newcommand{\AuditRawZero}{--}
\newcommand{\AuditRawNorm}{--}
\newcommand{\AuditPOVCos}{--}
\newcommand{\AuditPOVNeg}{--}
\newcommand{\AuditPOVZero}{--}
\newcommand{\AuditPOVNorm}{--}

% Prefix-horizon audit.  Conflict rates are percentages; cosine entries are
% batch/window aggregates computed against the independent outcome reference.
\newcommand{\HorizonTriggerRate}{--}
\newcommand{\HorizonMatchedPairs}{--}
\newcommand{\HorizonTriggeredConflictPre}{--}
\newcommand{\HorizonTriggeredConflictPost}{--}
\newcommand{\HorizonTriggeredConflictDelta}{--}
\newcommand{\HorizonTriggeredCosPre}{--}
\newcommand{\HorizonTriggeredCosPost}{--}
\newcommand{\HorizonControlConflictPre}{--}
\newcommand{\HorizonControlConflictPost}{--}
\newcommand{\HorizonControlConflictDelta}{--}
\newcommand{\HorizonControlCosPre}{--}
\newcommand{\HorizonControlCosPost}{--}
\newcommand{\HorizonConflictDID}{--}


We evaluate whether prefix--outcome validation corrects unreliable token-level OPD gradients and thereby improves on-policy distillation for reasoning. We compare \textbf{POV-OPD} with two complementary baselines: \textbf{GRPO}, which uses only sparse verifier feedback, and \textbf{OPD}, which uses dense teacher--student logit differences without outcome or prefix validation. We conduct experiments under four student--teacher configurations, spanning different model families and capacity gaps. Within each configuration, all methods use the same student initialization, training data, rollout budget, and optimization hyperparameters whenever applicable. Thus, the comparison isolates the contribution of validating dense teacher supervision.

\begin{table*}[t]
\centering
\small
\caption{Main results and compact mechanism audit. Panel (a) compares four student--teacher configurations; entries are percentages. Panels (b)--(d) use DeepSeek-R1-Distill-Qwen-1.5B with JustRL-1.5B. O/P/S denote outcome validation, prefix validation, and teacher-supported boosting. Panel (c) compares gradients on rollout set $A$ with an independent outcome-gradient reference from rollout set $B$ on the same prompts. ``Neg.'' and ``Zero'' denote negative-cosine and near-zero method-gradient rates, respectively.}
\label{tab:main-results}
\resizebox{\textwidth}{!}{%
\begin{tabular}{llcccccccc}
\hline
\textbf{Student / Teacher} & \textbf{Method} & \multicolumn{2}{c}{\textbf{AMC}} & \multicolumn{2}{c}{\textbf{AIME24}} & \multicolumn{2}{c}{\textbf{AIME25}} & \multicolumn{2}{c}{\textbf{GPQA}} \\
\cline{3-10}
& & Avg@8 & Best@8 & Avg@8 & Best@8 & Avg@8 & Best@8 & Avg@8 & Best@8 \\
\hline
Qwen3-1.7B / Qwen3-4B-T
& GRPO & -- & -- & -- & -- & -- & -- & -- & -- \\
& OPD & -- & -- & -- & -- & -- & -- & -- & -- \\
& \textbf{POV-OPD} & -- & -- & -- & -- & -- & -- & -- & -- \\
\hline
DS-R1-Qwen-1.5B / JustRL-1.5B
& GRPO & -- & -- & -- & -- & -- & -- & -- & -- \\
& OPD & -- & -- & -- & -- & -- & -- & -- & -- \\
& \textbf{POV-OPD} & -- & -- & -- & -- & -- & -- & -- & -- \\
\hline
DS-R1-Qwen-1.5B / Skywork-OR1-7B
& GRPO & -- & -- & -- & -- & -- & -- & -- & -- \\
& OPD & -- & -- & -- & -- & -- & -- & -- & -- \\
& \textbf{POV-OPD} & -- & -- & -- & -- & -- & -- & -- & -- \\
\hline
DS-R1-Qwen-7B / Skywork-OR1-7B
& GRPO & -- & -- & -- & -- & -- & -- & -- & -- \\
& OPD & -- & -- & -- & -- & -- & -- & -- & -- \\
& \textbf{POV-OPD} & -- & -- & -- & -- & -- & -- & -- & -- \\
\hline
\end{tabular}%
}
\vspace{4pt}

\begin{minipage}[t]{0.44\textwidth}
\centering
\textbf{(b) Component ablation}\\[-2pt]
\scriptsize
\setlength{\tabcolsep}{2.8pt}
\renewcommand{\arraystretch}{0.92}
\begin{tabular}{@{}lccccc@{}}
\hline
\textbf{Variant} & \textbf{O} & \textbf{P} & \textbf{S} & \textbf{Macro Avg@8} & \textbf{Macro Best@8} \\
\hline
OPD & -- & -- & -- & \AblOPDAvg & \AblOPDBest \\
$+$ Outcome & $\checkmark$ & -- & -- & \AblOutcomeAvg & \AblOutcomeBest \\
$+$ Prefix & -- & $\checkmark$ & -- & \AblPrefixAvg & \AblPrefixBest \\
$+$ Outcome $+$ Prefix & $\checkmark$ & $\checkmark$ & -- & \AblOutcomePrefixAvg & \AblOutcomePrefixBest \\
\textbf{POV-OPD} & $\checkmark$ & $\checkmark$ & $\checkmark$ & \AblFullAvg & \AblFullBest \\
\hline
\end{tabular}
\end{minipage}\hfill
\begin{minipage}[t]{0.54\textwidth}
\centering
\textbf{(c) Independent-reference gradient audit}\\[-2pt]
\scriptsize
\setlength{\tabcolsep}{3.0pt}
\renewcommand{\arraystretch}{0.92}
\begin{tabular}{@{}lcccc@{}}
\hline
\textbf{Signal on $A$} & \textbf{Cosine $\uparrow$} & \textbf{Neg. (\%) $\downarrow$} & \textbf{Zero (\%) $\downarrow$} & \textbf{Rel. norm} \\
\hline
Outcome-PG & $\AuditOutcomeCos$ & $\AuditOutcomeNeg$ & $\AuditOutcomeZero$ & $\AuditOutcomeNorm$ \\
Raw OPD & $\AuditRawCos$ & $\AuditRawNeg$ & $\AuditRawZero$ & $\AuditRawNorm$ \\
\textbf{POV-OPD} & $\AuditPOVCos$ & $\AuditPOVNeg$ & $\AuditPOVZero$ & $\AuditPOVNorm$ \\
\hline
\end{tabular}
\end{minipage}

\vspace{4pt}
\begin{minipage}[t]{0.94\textwidth}
\centering
\textbf{(d) Does $b^\star$ identify less outcome-compatible OPD gradients?}\\[-2pt]
\scriptsize
\setlength{\tabcolsep}{5.0pt}
\renewcommand{\arraystretch}{0.92}
\begin{tabular}{@{}lccccc@{}}
\hline
\textbf{Group} & \textbf{Conflict pre (\%)} & \textbf{Conflict post (\%)} & \textbf{$\Delta$ (pp)} & \textbf{Mean cos. pre} & \textbf{Mean cos. post} \\
\hline
Triggered & $\HorizonTriggeredConflictPre$ & $\HorizonTriggeredConflictPost$ & $\HorizonTriggeredConflictDelta$ & $\HorizonTriggeredCosPre$ & $\HorizonTriggeredCosPost$ \\
Position-matched control & $\HorizonControlConflictPre$ & $\HorizonControlConflictPost$ & $\HorizonControlConflictDelta$ & $\HorizonControlCosPre$ & $\HorizonControlCosPost$ \\
\hline
\end{tabular}
\end{minipage}
\end{table*}


\subsection{Experimental Setup}

\paragraph{Training.}
We train on DAPO-Math-17K \cite{yu2025dapo} and consider four configurations: (1) Qwen3-1.7B distilled from Qwen3-4B Thinking \cite{yang2025qwen3}; (2) DeepSeek-R1-Distill-Qwen-1.5B \cite{guo2025deepseekr1} distilled from JustRL-1.5B \cite{he2025justrl}; (3) DeepSeek-R1-Distill-Qwen-1.5B distilled from Skywork-OR1-Math-7B \cite{he2025skyworkor1}; and (4) DeepSeek-R1-Distill-Qwen-7B distilled from Skywork-OR1-Math-7B. For each prompt, the student samples on-policy responses and is evaluated by a rule-based mathematical verifier. GRPO uses the resulting outcome reward to construct group-relative advantages. OPD uses the teacher--student log-probability gap as a dense token-level signal. POV-OPD uses the same dense OPD signal, but outcome validation distinguishes outcome-aligned from anti-aligned teacher preferences and prefix validation attenuates supervision after sustained teacher-surprisal drift. In each configuration, GRPO uses the same student initialization but does not query the teacher; the listed teacher is used only by OPD and POV-OPD.

\paragraph{Evaluation.}
We evaluate on four reasoning benchmarks: AMC23 (AMC), AIME 2024 (AIME24), and AIME 2025 (AIME25) from the MAA competition series \cite{maa2025competitions}, and GPQA \cite{rein2023gpqa}. For each problem, we independently sample eight responses from the evaluated policy. \emph{Avg@8} is the mean correctness over the eight samples, whereas \emph{Best@8} reports whether at least one of the eight samples is correct, averaged over problems. Higher values are better for both metrics. We use the same decoding budget, prompt template, and rule-based answer verifier for all compared methods within each benchmark.

\subsection{Main Results}

Table~\ref{tab:main-results}(a) reports the main comparison across all four student--teacher configurations. The table jointly measures average reasoning reliability and the benefit of repeated sampling. Avg@8 measures the quality of a typical sampled response, while Best@8 measures the attainable performance when a small response budget is available. We report both because a method can improve the best sampled trajectory while degrading the average trajectory quality, which would be undesirable for practical reasoning systems.

\paragraph{Reading the comparison.}
The most direct evidence for POV-OPD is a consistent improvement over OPD on Avg@8: this would show that validating dense supervision improves the quality of typical on-policy trajectories rather than merely increasing sampling variance. Gains on Best@8 additionally indicate that the method preserves or improves exploration of successful solution paths. In contrast, a gain over GRPO demonstrates the value of dense teacher guidance beyond sparse outcome feedback. We further analyze how outcome and prefix validation reshape unreliable OPD gradients and isolate the contribution of each component in the following ablations.

\subsection{Component Ablation and Gradient Audit}
\label{sec:gradient-audit}

Panels~(b)--(d) of Table~\ref{tab:main-results} use DeepSeek-R1-Distill-Qwen-1.5B as the student and JustRL-1.5B as the teacher.  The component ablation changes only the validation weights.  Outcome-only disables prefix weighting and sets $\alpha=0$; prefix-only retains every raw OPD token, applies only $w^{\mathrm{pre}}$, and also sets $\alpha=0$; outcome+prefix keeps both validation weights but omits the support boost.  The final two rows therefore isolate the effect of $w^{\mathrm{sup}}$.  Macro Avg@8 and Best@8 give equal weight to each of the four benchmarks in the main evaluation.

\paragraph{Gradient protocol.}
At each pre-specified frozen student checkpoint, we use the same held-out prompts to sample two independent response groups, $A_j$ and $B_j$.  Method gradients are computed on $A_j$, while $B_j$ supplies an independent outcome-policy-gradient reference $g^{\mathrm{out}}_{j,B}$.  Outcome-PG on $A_j$ measures the agreement between two independent outcome estimators.  All advantages and validation weights are detached, and every objective uses the same trainable parameters, token masking, and loss normalization.  We first define one common audit set $\mathcal J=\{j:\lVert g^{\mathrm{out}}_{j,B}\rVert_2\geq\epsilon_{\mathrm{ref}}\}$ using only the reference norm, and use this set for every method.  For $X\in\{\mathrm{out},\mathrm{OPD},\mathrm{POV}\}$, let $\mathcal J_X\subseteq\mathcal J$ contain batches whose method gradient also exceeds a pre-specified numerical threshold.  We report
\begin{equation}
\begin{aligned}
c_j^X&=\frac{\langle g^X_{j,A},g^{\mathrm{out}}_{j,B}\rangle}
{\lVert g^X_{j,A}\rVert_2\lVert g^{\mathrm{out}}_{j,B}\rVert_2},\\
p_-^X&=\frac{1}{|\mathcal J_X|}\sum_{j\in\mathcal J_X}{\bf1}[c_j^X<0],
&z_0^X&=1-\frac{|\mathcal J_X|}{|\mathcal J|}.
\end{aligned}
\label{eq:gradient-audit}
\end{equation}
Cosine is averaged over $\mathcal J_X$ per audit batch rather than computed after averaging gradients.  Near-zero method gradients are not silently discarded: their rate $z_0^X$ is reported, while the common reference-valid count is $|\mathcal J|=\AuditN$.  We also report the batch-wise mean of $\lVert g^X_{j,A}\rVert/(\lVert g^{\mathrm{out}}_{j,B}\rVert+\epsilon)$ over all $j\in\mathcal J$, including method-gradient zeros, so that an apparent alignment gain cannot be attributed merely to shrinking the update toward zero.

\paragraph{Component effects.}
Outcome validation changes Macro Avg@8 from \AblOPDAvg{} to \AblOutcomeAvg{}, while prefix validation alone obtains \AblPrefixAvg{}.  Combining them yields \AblOutcomePrefixAvg{}, and adding teacher-supported boosting gives \AblFullAvg{}.  These comparisons respectively test outcome correction, prefix correction, their complementarity, and the incremental value of support; a component is treated as beneficial only when its gain exceeds the evaluation uncertainty.

\paragraph{Overall gradient correction.}
Against the independent outcome reference, raw OPD obtains cosine \(\AuditRawCos\) with \(\AuditRawNeg\%\) negative-cosine batches, whereas POV-OPD obtains \(\AuditPOVCos\) with \(\AuditPOVNeg\%\).  The corresponding relative norms are \(\AuditRawNorm\) and \(\AuditPOVNorm\).  This audit measures agreement with an outcome-based optimization direction, not access to a ground-truth token gradient; it is therefore interpreted as a mechanism check rather than a proof of token correctness.

\paragraph{Before and after the detected horizon.}
For a window $W_{i,b}$, define its raw OPD update as $g^{\mathrm{OPD}}_{i,b}=|W_{i,b}|^{-1}\sum_{t\in W_{i,b}}D_{i,t}\nabla_\theta\log\pi_\theta(y_{i,t}\mid h_{i,t})$.  It is conflicting iff $\langle g^{\mathrm{OPD}}_{i,b},g^{\mathrm{out}}_{j,B}\rangle<0$.  Panel (d) reports the equal-window mean of the corresponding cosines, separately before $b^\star$ and over the trigger window and all later windows.  Each triggered trajectory is matched to an untriggered rollout of the same prompt and outcome class with similar length; the control uses a pseudo-boundary at the same normalized position.  CUSUM triggers on \(\HorizonTriggerRate\%\) of audited trajectories, leaving \(\HorizonMatchedPairs\) matched pairs.  In triggered trajectories, the conflict rate changes from \(\HorizonTriggeredConflictPre\%\) before $b^\star$ to \(\HorizonTriggeredConflictPost\%\) after it, compared with \(\HorizonControlConflictPre\%\) to \(\HorizonControlConflictPost\%\) for the position-matched control.  The difference-in-differences is \(\HorizonConflictDID\) percentage points.  Evidence for localization requires this value to be positive with a 95\% confidence interval excluding zero and a larger post-horizon cosine deterioration than in the matched control.  Because prefix weights are positive constants within each window, they cannot alter that window's direction, conflict indicator, or cosine.  Panel (d) therefore tests whether $b^\star$ localizes raw OPD regions that are less outcome-compatible on average; the method attenuates the mass assigned to those post-horizon regions, while aggregate gradient correction is tested separately in panel (c).

